\documentclass[UTF8,fontset=fandol,zihao=-4]{ctexrep}

\usepackage[a4paper,left=2.7cm,right=2.7cm,top=2.6cm,bottom=2.6cm]{geometry}
\usepackage{setspace}
\usepackage{amsmath,amssymb,amsfonts}
\usepackage{booktabs,longtable,tabularx,array,multirow}
\usepackage[ruled,vlined,linesnumbered]{algorithm2e}
\usepackage{xcolor}
\usepackage{hyperref}
\usepackage{enumitem}
\usepackage{fancyhdr}
\usepackage{lastpage}
\usepackage{caption}
\usepackage{graphicx}

\hypersetup{colorlinks=true,linkcolor=black,urlcolor=black,citecolor=black}
\onehalfspacing
\setlength{\parindent}{2em}
\setlength{\parskip}{0pt}
\setlength{\tabcolsep}{3pt}
\setlength{\headheight}{15pt}
\setlist{nosep,leftmargin=2em}

\ctexset{
  chapter = {
    format = \centering\zihao{3}\bfseries,
    name = {第,章},
    number = \chinese{chapter},
    aftername = \quad,
    beforeskip = 18pt,
    afterskip = 24pt
  },
  section = {
    format = \zihao{4}\bfseries,
    beforeskip = 12pt,
    afterskip = 8pt
  },
  subsection = {
    format = \zihao{-4}\bfseries,
    beforeskip = 10pt,
    afterskip = 6pt
  }
}

\pagestyle{fancy}
\fancyhf{}
\fancyhead[L]{\small 链甲 AgentArmor}
\fancyhead[C]{\small Multi-Agent Shared Write}
\fancyhead[R]{\small 阶段性报告}
\fancyfoot[C]{\small 第 \thepage 页 \quad 共 \pageref{LastPage} 页}
\renewcommand{\headrulewidth}{0.4pt}
\renewcommand{\footrulewidth}{0pt}

\SetAlgorithmName{Algorithm}{Algorithm}{算法目录}
\SetAlgoNlRelativeSize{0}
\SetNlSty{textbf}{}{}
\SetAlFnt{\small}
\SetKwInOut{KwIn}{Input}
\SetKwInOut{KwOut}{Output}
\SetKw{KwReturn}{return}
\SetKw{KwBreak}{break}
\SetKwComment{tcp}{$\triangleright$ }{}

\newcommand{\sys}{\textsc{MASW}}
\newcommand{\attack}{\textsc{MASW-Poison-Exec}}
\newcommand{\tool}[1]{\texttt{#1}}
\newcommand{\vect}[1]{\boldsymbol{#1}}

\begin{document}

\begin{titlepage}
  \centering
  \vspace*{3.5cm}
  {\zihao{1}\bfseries 链甲 AgentArmor\par}
  \vspace{0.9cm}
  {\zihao{2}\bfseries Multi-Agent Shared Write 安全防御\par}
  \vspace{0.8cm}
  {\zihao{3}\bfseries 阶段性设计与测试报告\par}
  \vfill
  {\zihao{4} 2026年6月28日\par}
\end{titlepage}

\pagenumbering{Roman}
\tableofcontents
\clearpage
\pagenumbering{arabic}

\chapter*{摘要}
\addcontentsline{toc}{chapter}{摘要}

多 Agent 系统正在从单轮问答扩展到长期协作、共享记忆和自动工具调用。与传统单 Agent 场景相比，多个 Agent 共同读写同一份记忆时，会产生一种新的信任传播风险：低信任外部内容经过某个 Agent 的读取、摘要或改写后，被写入共享记忆；后续 Agent 再将这段记忆当作可信事实或执行依据，从而触发越权工具调用、数据外传或长期污染。

本文针对该类风险定义 Type 5: Multi-Agent Shared Write 场景，并提出 \sys{} 防御框架。框架围绕“不可信内容不得因 Agent 改写或写入而自动提升为可信记忆”这一安全不变量，将防御链拆分为输入标记、候选事实抽取、风险过滤、溯源写入网关、检索审计、执行仲裁和事后撤销七个阶段。与单纯文本分类器不同，\sys{} 的核心目标不是只判断某句话是否危险，而是在共享记忆的写入、检索和执行三个关键边界上阻断攻击链。

为了验证系统有效性，本文构建了最小可用测试集，包含 prompt injection、tool misuse、memory poisoning、agent hijacking 四类攻击样本各 20 条，并加入 45 条良性样本；在技术选型和基线对比中进一步扩展到 145 条样本。实验包括单元测试、消融实验、节点效果分析、技术选型、基线对比、延迟测试和 8 个端到端攻击场景。结果显示，\sys{} 在当前数据集上阻断了全部攻击样本的共享写入；在端到端验证中，脆弱系统发生 8/8 次危险执行，而 \sys{} 将危险执行降为 0/8。

\textbf{关键词：} 多 Agent 安全；共享记忆；提示注入；记忆投毒；工具调用；溯源验证

\chapter{介绍}

\section{研究背景}

多 Agent 系统通常由多个具备不同任务角色的 Agent 共同完成复杂流程。例如，一个 Agent 负责检索网页，一个 Agent 负责总结知识，一个 Agent 负责生成执行计划，另一个 Agent 负责调用邮件、数据库或代码仓库工具。为了提升协作效率，这类系统往往会引入共享对话日志、crew memory、共享知识库或 RAG memory store，使不同 Agent 可以复用彼此产生的中间结果。

共享记忆提升了效率，也放大了安全风险。单 Agent 场景中的提示注入通常只影响当前对话；而在 Multi-Agent Shared Write 场景中，一次恶意输入可能被写入长期共享状态，并在未来的正常任务中被其他 Agent 召回。这使攻击从“即时影响一次回答”升级为“跨 Agent、跨任务、可持久化触发”的系统风险。

\section{问题定义}

本文研究的系统流程如下：

\[
\mathrm{External\ Data}
\rightarrow
\mathrm{Agent\ A}
\rightarrow
\mathrm{Shared\ Memory}
\rightarrow
\mathrm{Agent\ B/C}
\rightarrow
\mathrm{Tool\ Action}
\]

在脆弱系统中，Agent A 只要拥有写入权限，其输出就可能直接进入共享记忆；Agent B/C 在后续任务中检索到该记忆后，又可能把它当成可信指令或事实使用。攻击者不需要直接调用高权限工具，只需要控制网页、工具返回、邮件或文档中的一段文本，即可通过低权限 Agent 完成间接影响。

因此，本文关注的核心问题是：

\begin{quote}
在多个 Agent 平等写入共享记忆的系统中，如何防止低信任外部内容被自动提升为高信任共享记忆，并进一步诱导后续 Agent 执行危险操作？
\end{quote}

\section{主要贡献}

本文当前阶段的主要贡献包括：

\begin{enumerate}
  \item 提出 Multi-Agent Shared Write 场景下的信任提升风险模型，明确指出共享记忆的危险不在于“存储文本”，而在于“文本被后续 Agent 当作可信上下文使用”。
  \item 构造 \attack{} 端到端攻击链，复现外部注入、共享记忆污染、后续检索和危险工具执行的完整流程。
  \item 设计 \sys{} 分层防御框架，在写入、检索和执行三个边界分别设置溯源验证、信任审计和动作仲裁。
  \item 构建最小可用评测数据集和端到端验证场景，使用消融实验、节点效果、基线对比和延迟测试验证当前设计。
\end{enumerate}

\chapter{背景与威胁模型}

\section{Multi-Agent Shared Write 场景}

Multi-Agent Shared Write 指多个 Agent 均可向共享记忆写入内容，且写入内容会在后续任务中被其他 Agent 检索和使用。其典型特征包括：

\begin{enumerate}
  \item \textbf{写入主体多样：} 任意 Agent 都可能成为共享记忆的写入者。
  \item \textbf{内容来源复杂：} 写入内容可能来自用户、网页、搜索结果、工具返回、文件或模型总结。
  \item \textbf{使用时机延迟：} 污染内容不一定立即触发危害，可能在未来任务中被召回。
  \item \textbf{权限不对称：} 写入内容的 Agent 权限可能较低，读取内容并执行工具的 Agent 权限可能较高。
\end{enumerate}

这类系统的隐含假设是：共享记忆中的内容具有较高可信度。但该假设在外部内容可被写入时不成立。

\section{攻击者能力}

本文假设攻击者可以控制至少一个外部内容源，例如网页、搜索结果、工具返回、文档片段、邮件正文或 OCR 文本。攻击者能够诱导低权限 Agent A 读取该内容，并构造自然语言、中英文混合、伪装业务规则、伪装知识库同步流程等 payload。

攻击者不能直接修改系统策略，不能直接调用高权限工具，也不能修改 \sys{} 防御代码。换言之，本文研究的不是系统被完全攻破后的情况，而是外部不可信内容经由正常 Agent 流程进入共享记忆后造成的间接攻击。

\section{攻击目标}

攻击者目标可以分为四类：

\begin{enumerate}
  \item \textbf{完整性破坏：} 将恶意规则、伪政策或伪任务规范写入共享记忆。
  \item \textbf{机密性破坏：} 诱导后续 Agent 外传上下文、摘要、密钥或内部数据。
  \item \textbf{权限绕过：} 让高权限 Agent 调用 \tool{email.send}、\tool{database.write}、\tool{repo.commit}、\tool{deploy.production}、\tool{secret.read} 等工具。
  \item \textbf{持久化影响：} 使污染内容跨任务、跨 Agent 留存，并在未来被正常查询触发。
\end{enumerate}

\section{攻击面分析}

\begin{longtable}{@{}p{0.17\textwidth}p{0.21\textwidth}p{0.25\textwidth}p{0.24\textwidth}@{}}
\caption{Multi-Agent Shared Write 攻击面分析}\\
\toprule
\textbf{阶段} & \textbf{系统状态} & \textbf{潜在威胁} & \textbf{失败根因}\\
\midrule
\endfirsthead
\toprule
\textbf{阶段} & \textbf{系统状态} & \textbf{潜在威胁} & \textbf{失败根因}\\
\midrule
\endhead
外部输入 & 网页、文件、工具返回、邮件 & 隐藏指令、伪装规范、跨语言注入 & 内容缺少来源、信任和污点标签\\
Agent 处理 & 摘要、抽取、改写 & 指令被改写成事实 & 生成文本被误当作验证事实\\
共享写入 & 候选内容进入 shared memory & 自动信任提升、污染持久化 & 写入权限扁平，缺少网关\\
记忆检索 & 后续 Agent 召回记忆 & 污染内容被当成上下文 & 检索只看相关性，不看风险\\
工具执行 & Agent 生成 action proposal & 数据外传、越权写入、部署操作 & Agent 输出被直接当作授权\\
事后处置 & 污染已传播 & 派生内容无法清理 & 缺少溯源图和影响分析\\
\bottomrule
\end{longtable}

\section{安全不变量}

\sys{} 的核心安全不变量为：

\[
\textbf{No Automatic Trust Promotion}
\]

即：不可信外部内容不能仅因为经过 Agent 摘要、改写或写入，就自动变成可信共享记忆。

设内容项为 \(x\)，来源信任为 \(T_{\mathrm{src}}(x)\)，写入者信任为 \(T_{\mathrm{writer}}(a)\)，共享记忆信任为 \(T_{\mathrm{mem}}(x)\)。脆弱系统常见的错误规则为：

\[
T_{\mathrm{mem}}(x)=T_{\mathrm{writer}}(a)
\]

该规则忽略了内容来源，使得低信任外部数据可以借助高权限或可写 Agent 完成信任漂白。本文将其改写为：

\[
T_{\mathrm{mem}}(x)=
\begin{cases}
\mathrm{promote}(x), & V(x)=1 \land R(x)\le \tau_w,\\
T_{\mathrm{src}}(x), & \text{otherwise},
\end{cases}
\]

其中 \(V(x)\) 表示溯源验证是否通过，\(R(x)\) 表示风险分数，\(\tau_w\) 表示共享写入阈值。

\chapter{系统设计}

\section{设计目标}

\sys{} 的设计目标不是替代所有 Agent 框架，也不是对所有文本做统一拦截，而是在 Multi-Agent Shared Write 场景中建立清晰的安全边界。系统需要满足以下要求：

\begin{enumerate}
  \item 外部内容进入系统时必须带有来源、信任和污点标签。
  \item Agent 生成的中间结果只能作为候选事实，不能直接成为共享可信记忆。
  \item 共享记忆写入必须经过独立验证，不能由写入 Agent 自行决定可信度。
  \item 检索结果只能作为证据，不能天然作为执行指令。
  \item 高影响工具调用必须经过动作仲裁和权限检查。
  \item 污染确认后，系统能够追踪并撤销派生记忆。
\end{enumerate}

\section{总体架构}

\sys{} 将完整防御链划分为七个阶段：

\begin{enumerate}
  \item \textbf{D1 输入标记：} 为外部输入绑定 source、channel、trust、taint 和 scope。
  \item \textbf{D2 候选事实抽取：} 将 Agent A 的输出限制为候选事实，不允许直接写入共享记忆。
  \item \textbf{D3 风险过滤：} 检测提示注入、工具滥用、外传、记忆投毒和 Agent 劫持信号。
  \item \textbf{D4 溯源写入网关：} 依据证据、来源、冲突和风险决定是否允许写入。
  \item \textbf{D5 检索审计：} 检索时同时考虑相关性、信任等级、作用域和风险。
  \item \textbf{D6 执行仲裁：} 高影响动作必须与用户目标、权限和上下文一致。
  \item \textbf{D7 撤销与影响分析：} 基于溯源图追踪污染节点及其派生内容。
\end{enumerate}

这七个阶段对应攻击链中的不同断点：D1--D4 控制污染是否进入共享记忆，D5--D6 控制污染是否能影响执行，D7 控制污染发生后的扩散影响。

\section{核心数据抽象}

为了避免自然语言内容在不同阶段被误用，\sys{} 将系统状态分为四类对象：

\[
X=\langle text,source,channel,trust,taint,scope\rangle
\]

\[
C=\langle claim,evidence,parent,taint,risk\_hint\rangle
\]

\[
M=\langle claim,provenance,trust,risk,scope,parent\_ids\rangle
\]

\[
A=\langle tool,args,user\_goal,context,impact\rangle
\]

其中 \(X\) 表示带标签输入，\(C\) 表示候选事实，\(M\) 表示共享记忆记录，\(A\) 表示工具动作提案。该抽象的作用是把“文本内容”和“可信状态”拆开：同一段文本在不同对象中具有不同权限，只有通过 D4 的候选事实才能转化为共享记忆。

\section{写入、检索与执行三重边界}

\sys{} 的关键不是某个单点检测器，而是三重边界：

\begin{enumerate}
  \item \textbf{写入边界：} 外部内容无法绕过 D3/D4 直接进入共享记忆。
  \item \textbf{检索边界：} 低信任或跨作用域记忆无法被当作高信任上下文使用。
  \item \textbf{执行边界：} 检索到的记忆不能自动授权高影响工具调用。
\end{enumerate}

因此，即使某一层被绕过，后续层仍可阻断攻击链。

\chapter{攻击分类与防御算法}

\section{符号说明}

\begin{longtable}{@{}p{0.14\textwidth}p{0.30\textwidth}p{0.14\textwidth}p{0.29\textwidth}@{}}
\caption{伪代码符号表}\\
\toprule
\textbf{符号} & \textbf{含义} & \textbf{符号} & \textbf{含义}\\
\midrule
\endfirsthead
\toprule
\textbf{符号} & \textbf{含义} & \textbf{符号} & \textbf{含义}\\
\midrule
\endhead
\(\mathcal{X}\) & 带标签外部输入集合 & \(\mathcal{C}\) & 候选事实集合\\
\(\mathcal{M}\) & 共享记忆集合 & \(\mathcal{Q}\) & 隔离队列\\
\(\mathcal{A}\) & 动作提案集合 & \(\mathcal{G}\) & 记忆溯源图\\
\(T(x)\) & 输入或记忆的信任分数 & \(R(c)\) & 候选事实风险分数\\
\(\tau_3\) & D3 风险过滤阈值 & \(\tau_4\) & D4 写入阈值\\
\(E(c)\) & 证据验证谓词 & \(S(c)\) & 来源验证谓词\\
\(K(c)\) & 冲突检测谓词 & \(P(a)\) & 动作权限谓词\\
\bottomrule
\end{longtable}

\section{端到端攻击流程}

\attack{} 攻击的目标是让恶意外部内容跨过 Agent A，进入共享记忆，并在 Agent B 的正常任务中被召回后触发危险工具。该攻击过程可形式化为 Algorithm~\ref{alg:attack}。

\begin{algorithm}[H]
\caption{\attack{} 攻击复现流程}
\label{alg:attack}
\KwIn{外部内容源 \(d\)，注入载荷 \(p\)，低权限 Agent \(A_1\)，高权限 Agent \(A_2\)，共享记忆 \(\mathcal{M}\)，目标查询 \(q_t\)}
\KwOut{危险动作 \(a^\star\) 或攻击失败标记 \(\bot\)}
\BlankLine
\(d' \leftarrow \mathrm{Inject}(d,p)\)\tcp*{攻击者把伪规则写入网页、工具返回或文档}
\(s \leftarrow A_1.\mathrm{ReadAndSummarize}(d')\)\tcp*{低权限 Agent 将载荷改写为经验或规范}
\If{\(\mathrm{VulnerableWritePolicy}(s)=\mathrm{true}\)}{
  \(\mathcal{M}\leftarrow \mathcal{M}\cup\{s\}\)\tcp*{污染内容进入共享记忆}
}
\(\mathcal{R}\leftarrow A_2.\mathrm{Retrieve}(\mathcal{M},q_t)\)\;
\(a^\star \leftarrow A_2.\mathrm{PlanAction}(q_t,\mathcal{R})\)\;
\If{\(a^\star.tool\in\{\tool{email.send},\tool{database.write},\tool{deploy.production},\tool{secret.read}\}\)}{
  \KwReturn \(a^\star\)\tcp*{攻击成功：高权限工具被污染记忆诱导}
}
\KwReturn \(\bot\)\;
\end{algorithm}

该攻击能够成立，是因为脆弱系统同时满足两个条件：一是共享写入阶段缺少信任提升验证；二是工具执行阶段把检索到的记忆当成授权依据。

\section{共享写入防御算法}

共享写入是整条攻击链中最关键的边界。Algorithm~\ref{alg:write} 展示了 \sys{} 的写入流程：外部内容先被标记，再被抽取为候选事实，随后经过风险过滤和溯源网关，只有通过验证的内容才能进入共享记忆。

\begin{algorithm}[H]
\caption{Provenance-Gated Shared Write}
\label{alg:write}
\KwIn{外部输入 \(x\)，写入 Agent \(a\)，共享记忆 \(\mathcal{M}\)，隔离队列 \(\mathcal{Q}\)}
\KwOut{更新后的 \(\mathcal{M}\) 与 \(\mathcal{Q}\)}
\BlankLine
\(x' \leftarrow \mathrm{Label}(x)\)\tcp*{D1：绑定 source、trust、taint、scope}
\(\mathcal{C}\leftarrow \mathrm{ExtractCandidateFacts}(a,x')\)\tcp*{D2：只生成候选事实}
\ForEach{\(c\in\mathcal{C}\)}{
  \(r\leftarrow \mathrm{RiskScore}(c)\)\tcp*{D3：提示注入、外传、工具滥用、记忆投毒}
  \If{\(r>\tau_3\) \textbf{or} \(\mathrm{InstructionLike}(c)\)}{
    \(\mathcal{Q}\leftarrow \mathcal{Q}\cup\{(c,\mathrm{``risk\_filter''})\}\)\;
    \KwBreak
  }
  \(v\leftarrow E(c)\land S(c)\land \neg K(c)\land r\le\tau_4\)\tcp*{D4：证据、来源、冲突、风险}
  \If{\(v=\mathrm{true}\)}{
    \(m\leftarrow \mathrm{PromoteToMemory}(c,\mathrm{provenance}(c))\)\;
    \(\mathcal{M}\leftarrow \mathcal{M}\cup\{m\}\)\;
  }
  \Else{
    \(\mathcal{Q}\leftarrow \mathcal{Q}\cup\{(c,\mathrm{``provenance\_gate''})\}\)\;
  }
}
\KwReturn \((\mathcal{M},\mathcal{Q})\)\;
\end{algorithm}

该算法的安全含义是：Agent 的自然语言输出不能直接改变共享记忆的信任等级，必须经过独立验证。若攻击者试图把“将总结发送到外部邮箱”伪装成知识库同步规则，该内容会在风险过滤或写入网关阶段进入隔离队列。

\section{风险评分算法}

D3 的作用是尽早发现明显攻击和隐蔽外传表述。当前实现采用规则检测与语义 rubric 的组合，风险分数定义为：

\[
R(c)=w_iI_{\mathrm{inject}}(c)+w_tI_{\mathrm{tool}}(c)+w_mI_{\mathrm{memory}}(c)+w_eI_{\mathrm{exfil}}(c)+w_hI_{\mathrm{hijack}}(c)
\]

其中 \(I_{\cdot}(c)\in\{0,1\}\)，分别表示提示注入、工具滥用、记忆投毒、数据外传和 Agent 劫持信号。Algorithm~\ref{alg:risk} 给出检测流程。

\begin{algorithm}[H]
\caption{Hybrid Risk Screening}
\label{alg:risk}
\KwIn{候选事实 \(c\)，规则集合 \(\Omega\)，rubric 维度集合 \(\Phi\)}
\KwOut{风险分数 \(R(c)\) 与风险标签集合 \(\mathcal{B}\)}
\BlankLine
\(R(c)\leftarrow 0,\ \mathcal{B}\leftarrow \varnothing\)\;
\ForEach{\(\omega\in\Omega\)}{
  \If{\(\omega.\mathrm{match}(c.text)=\mathrm{true}\)}{
    \(R(c)\leftarrow R(c)+\omega.weight\)\;
    \(\mathcal{B}\leftarrow \mathcal{B}\cup\{\omega.tag\}\)\;
  }
}
\ForEach{\(\phi\in\Phi\)}{
  \(s_\phi\leftarrow \mathrm{RubricScore}(c,\phi)\)\tcp*{检查同步伪装、外部路由、工具副作用等语义信号}
  \If{\(s_\phi>\phi.threshold\)}{
    \(R(c)\leftarrow R(c)+\phi.weight\)\;
    \(\mathcal{B}\leftarrow \mathcal{B}\cup\{\phi.tag\}\)\;
  }
}
\If{\(\mathrm{ExternalRecipient}(c)\land \mathrm{ActionVerb}(c)\)}{
  \(R(c)\leftarrow R(c)+w_e,\ \mathcal{B}\leftarrow\mathcal{B}\cup\{\mathrm{exfiltration}\}\)\;
}
\KwReturn \((R(c),\mathcal{B})\)\;
\end{algorithm}

基础实验显示，纯规则检测对显式攻击足够有效，但在“知识库同步”“团队归档”“自动推送”等隐蔽外传表达上存在漏报。因此 D3 使用 hybrid 策略；D4 写入网关则保持更保守的规则检测，以减少良性内部同步被误伤。

\section{检索审计算法}

即使污染内容绕过写入阶段，也不应在检索阶段被无条件交给后续 Agent。D5 将检索排序由单一相关性扩展为信任感知评分：

\[
S(m,q)=\mathrm{sim}(m,q)+\alpha T(m)-\beta R(m)-\gamma I_{\mathrm{scope}}(m,q)
\]

其中 \(\mathrm{sim}(m,q)\) 是相关性，\(T(m)\) 是记忆信任，\(R(m)\) 是风险分数，\(I_{\mathrm{scope}}\) 表示作用域不匹配。

\begin{algorithm}[H]
\caption{Trust-Aware Retrieval Audit}
\label{alg:retrieval}
\KwIn{目标查询 \(q\)，共享记忆 \(\mathcal{M}\)，最小信任阈值 \(\tau_T\)，最大风险阈值 \(\tau_R\)}
\KwOut{可进入上下文的记忆集合 \(\mathcal{R}\)}
\BlankLine
\(\mathcal{R}\leftarrow \varnothing\)\;
\ForEach{\(m\in\mathcal{M}\)}{
  \If{\(\mathrm{ScopeMismatch}(m,q)=\mathrm{true}\)}{
    \textbf{continue}\tcp*{跨任务或跨租户内容不进入上下文}
  }
  \If{\(T(m)<\tau_T\) \textbf{or} \(R(m)>\tau_R\)}{
    \(\mathrm{AuditLog}(m,q,\mathrm{``filtered''})\)\;
    \textbf{continue}\;
  }
  \(score\leftarrow \mathrm{sim}(m,q)+\alpha T(m)-\beta R(m)\)\;
  \(\mathcal{R}\leftarrow \mathcal{R}\cup\{(m,score,\mathrm{evidence\_only})\}\)\;
}
\(\mathcal{R}\leftarrow \mathrm{TopK}(\mathcal{R})\)\;
\KwReturn \(\mathcal{R}\)\;
\end{algorithm}

该算法保证检索结果的身份是 evidence，而不是 instruction。后续 Agent 可以参考其事实内容，但不能把它当作工具授权规则。

\section{执行仲裁算法}

D6 处理的是混淆代理风险。即使某条污染记忆进入上下文，只要它诱导高影响工具调用，动作仍必须通过独立仲裁。设动作 \(a\) 的允许条件为：

\[
\mathrm{Allow}(a)=U(a)\land P(a)\land \neg H(a)\land \neg X(a)
\]

其中 \(U(a)\) 表示与用户目标一致，\(P(a)\) 表示权限满足，\(H(a)\) 表示高影响动作，\(X(a)\) 表示外传或危险副作用。

\begin{algorithm}[H]
\caption{Execution Mediation for High-Impact Tools}
\label{alg:execution}
\KwIn{动作提案 \(a\)，用户目标 \(g\)，检索上下文 \(\mathcal{R}\)，工具策略 \(\Pi\)}
\KwOut{执行决策 \(d\in\{\mathrm{allow},\mathrm{deny},\mathrm{review}\}\)}
\BlankLine
\If{\(\neg\mathrm{GoalAligned}(a,g)\)}{
  \KwReturn \(\mathrm{deny}\)\tcp*{动作与用户任务不一致}
}
\If{\(\neg\mathrm{PermissionOK}(a,\Pi)\)}{
  \KwReturn \(\mathrm{deny}\)\tcp*{权限不足}
}
\If{\(\mathrm{ExternalSink}(a)\land \mathrm{SensitiveContext}(\mathcal{R})\)}{
  \KwReturn \(\mathrm{review}\)\tcp*{存在数据外传风险}
}
\If{\(a.tool\in\{\tool{email.send},\tool{database.write},\tool{repo.commit},\tool{deploy.production},\tool{secret.read}\}\)}{
  \KwReturn \(\mathrm{review}\)\tcp*{高影响工具必须人工确认}
}
\KwReturn \(\mathrm{allow}\)\;
\end{algorithm}

该算法并不依赖 D3/D4 是否成功。只要动作具有高影响或外传风险，就不会因为某段共享记忆声称“这是团队规范”而自动执行。

\section{撤销与影响分析算法}

D7 将共享记忆视为有向图 \(\mathcal{G}=(V,E)\)，边 \(u\rightarrow v\) 表示记忆 \(v\) 由 \(u\) 派生。污染节点 \(p\) 的影响集合为：

\[
\mathrm{Impact}(p)=\{v\in V\mid p\leadsto v\}
\]

\begin{algorithm}[H]
\caption{Revocation and Impact Analysis}
\label{alg:revocation}
\KwIn{溯源图 \(\mathcal{G}\)，污染记忆节点 \(p\)，共享记忆 \(\mathcal{M}\)}
\KwOut{撤销后的共享记忆 \(\mathcal{M}'\) 与影响集合 \(\mathcal{I}\)}
\BlankLine
\(\mathcal{I}\leftarrow \{p\},\ W\leftarrow \{p\}\)\;
\While{\(W\neq\varnothing\)}{
  \(u\leftarrow \mathrm{pop}(W)\)\;
  \ForEach{\(v\in\mathrm{Children}_{\mathcal{G}}(u)\)}{
    \If{\(v\notin\mathcal{I}\)}{
      \(\mathcal{I}\leftarrow\mathcal{I}\cup\{v\},\ W\leftarrow W\cup\{v\}\)\;
    }
  }
}
\ForEach{\(m\in\mathcal{I}\)}{
  \(\mathcal{M}\leftarrow \mathcal{M}\setminus\{m\}\)\;
  \(\mathrm{AuditLog}(m,\mathrm{``revoked''})\)\;
}
\KwReturn \((\mathcal{M},\mathcal{I})\)\;
\end{algorithm}

D7 不阻止首次攻击，但解决污染发现后的持续传播问题。没有 parent ids 和 provenance graph 时，系统只能删除单条记忆，无法判断哪些后续内容也被污染。

\chapter{实验与评估}

\section{研究问题}

本文通过六个研究问题评估 \sys{}：

\begin{itemize}
  \item \textbf{RQ1：} 每个防御节点能否防住对应攻击阶段？
  \item \textbf{RQ2：} 系统能否阻断共享记忆污染？
  \item \textbf{RQ3：} 组合防御是否优于单点防御？
  \item \textbf{RQ4：} D3/D4 的检测技术路线应如何选择？
  \item \textbf{RQ5：} 与同类工具基线相比，\sys{} 是否更有效？
  \item \textbf{RQ6：} 真实两轮攻击链能否被复现并阻断？
\end{itemize}

\section{数据集构造}

当前阶段先从最小可用规模开始构建数据集。基础集共 125 条，其中攻击样本 80 条，良性样本 45 条。技术选型和基线对比额外加入 20 条隐蔽同步、语义改写和边界良性样本，形成 145 条 hard set。

\begin{longtable}{@{}p{0.26\textwidth}p{0.14\textwidth}p{0.48\textwidth}@{}}
\caption{数据集类别与覆盖内容}\\
\toprule
\textbf{类别} & \textbf{数量} & \textbf{覆盖内容}\\
\midrule
\endfirsthead
\toprule
\textbf{类别} & \textbf{数量} & \textbf{覆盖内容}\\
\midrule
\endhead
Prompt Injection & 20 & 忽略原指令、伪系统消息、隐藏规则、跨语言注入、上下文泄露。\\
Tool Misuse & 20 & 邮件外发、数据库写入、仓库提交、生产部署、密钥读取。\\
Memory Poisoning & 20 & 伪知识库规则、伪同步策略、覆盖安全策略、长期持久化规范。\\
Agent Hijacking & 20 & 改变后续 Agent 目标、转发上下文、执行攻击者任务。\\
Benign & 45 & 正常搜索总结、知识整理、内部同步、无危险副作用任务。\\
\bottomrule
\end{longtable}

公开数据集如 PromptInject 和 JailbreakBench 可提供提示注入表达参考，但它们通常不包含共享记忆写入、检索、隔离、工具执行等标签。因此本项目采用“公开攻击模式参考 + 本系统场景重构”的方式构造样本。

\section{评价指标}

本文使用 Precision、Recall、F1 和 FPR 衡量拦截效果：

\[
\mathrm{Precision}=\frac{TP}{TP+FP},\quad
\mathrm{Recall}=\frac{TP}{TP+FN}
\]

\[
\mathrm{F1}=\frac{2\cdot\mathrm{Precision}\cdot\mathrm{Recall}}
{\mathrm{Precision}+\mathrm{Recall}},\quad
\mathrm{FPR}=\frac{FP}{FP+TN}
\]

端到端实验额外记录 poisoned memory write 和 dangerous execution；性能实验记录 P50、P95、P99 延迟。

\section{RQ1：节点单独效果}

\begin{longtable}{@{}p{0.18\textwidth}p{0.15\textwidth}p{0.085\textwidth}p{0.085\textwidth}p{0.085\textwidth}p{0.085\textwidth}p{0.17\textwidth}@{}}
\caption{节点单独效果}\\
\toprule
\textbf{节点配置} & \textbf{对应阶段} & \textbf{Prec.} & \textbf{Rec.} & \textbf{F1} & \textbf{FPR} & \textbf{结论}\\
\midrule
\endfirsthead
\toprule
\textbf{节点配置} & \textbf{对应阶段} & \textbf{Prec.} & \textbf{Rec.} & \textbf{F1} & \textbf{FPR} & \textbf{结论}\\
\midrule
\endhead
D1 only & 输入标记 & 0.6400 & 1.0000 & 0.7805 & 1.0000 & 不能独立拦截\\
D2 only & 候选抽取 & 0.0000 & 0.0000 & 0.0000 & 0.0000 & 中间层\\
D3 only & 风险过滤 & 1.0000 & 1.0000 & 1.0000 & 0.0000 & 基础集有效\\
D4 gate & 写入网关 & 1.0000 & 1.0000 & 1.0000 & 0.0000 & 核心边界\\
D5 only & 检索审计 & 1.0000 & 0.0625 & 0.1176 & 0.0000 & 后置防线\\
D6 only & 执行仲裁 & 0.6400 & 1.0000 & 0.7805 & 1.0000 & 单独误伤高\\
D7 & 撤销处置 & 1.0000 & 1.0000 & 1.0000 & 0.0000 & 事故后收束\\
\bottomrule
\end{longtable}

结果表明，D3 和 D4 是当前最直接的拦截边界；D1、D2 是前置结构化条件；D5、D6、D7 更偏向纵深防御和事故后收束，不能单独替代写入网关。

\section{RQ2 与 RQ3：消融实验}

\begin{longtable}{@{}p{0.23\textwidth}p{0.19\textwidth}p{0.085\textwidth}p{0.085\textwidth}p{0.085\textwidth}p{0.085\textwidth}@{}}
\caption{Config-1 到 Config-5 消融实验}\\
\toprule
\textbf{Config} & \textbf{启用阶段} & \textbf{Prec.} & \textbf{Rec.} & \textbf{F1} & \textbf{FPR}\\
\midrule
\endfirsthead
\toprule
\textbf{Config} & \textbf{启用阶段} & \textbf{Prec.} & \textbf{Rec.} & \textbf{F1} & \textbf{FPR}\\
\midrule
\endhead
Config-1 Baseline & none & 0.0000 & 0.0000 & 0.0000 & 0.0000\\
Config-2 D3 only & D2+D3 & 1.0000 & 1.0000 & 1.0000 & 0.0000\\
Config-3 D4 gate & D1+D2+D4 & 1.0000 & 1.0000 & 1.0000 & 0.0000\\
Config-4 D3+D4 & D1--D4 & 1.0000 & 1.0000 & 1.0000 & 0.0000\\
Config-5 ALL & D1--D7 & 1.0000 & 1.0000 & 1.0000 & 0.0000\\
\bottomrule
\end{longtable}

基础集上 D3/D4 已经饱和，因此严格意义上的 \(F1_{\mathrm{ALL}}>F1_{\mathrm{best}}\) 不成立。该结果不能被夸大为组合防御在所有指标上都具有数学协同，而应解释为：基础集难度不足以体现后置边界的贡献。端到端攻击链验证显示，当写入、检索和执行串联起来时，组合防御仍是必要的。

\section{RQ4：技术选型}

\begin{longtable}{@{}p{0.28\textwidth}p{0.10\textwidth}p{0.10\textwidth}p{0.10\textwidth}p{0.10\textwidth}p{0.17\textwidth}@{}}
\caption{检测器技术选型结果}\\
\toprule
\textbf{Config} & \textbf{Prec.} & \textbf{Rec.} & \textbf{F1} & \textbf{FPR} & \textbf{结论}\\
\midrule
\endfirsthead
\toprule
\textbf{Config} & \textbf{Prec.} & \textbf{Rec.} & \textbf{F1} & \textbf{FPR} & \textbf{结论}\\
\midrule
\endhead
D3\_RULE & 1.0000 & 0.8889 & 0.9412 & 0.0000 & 漏掉隐蔽表达\\
D3\_HYBRID & 1.0000 & 1.0000 & 1.0000 & 0.0000 & 推荐用于 D3\\
D4\_GATE\_RULE & 1.0000 & 1.0000 & 1.0000 & 0.0000 & 推荐用于 D4\\
D4\_GATE\_HYBRID & 0.9783 & 1.0000 & 0.9890 & 0.0364 & 写入误伤升高\\
D3 Hybrid + D4 Rule & 1.0000 & 1.0000 & 1.0000 & 0.0000 & 当前推荐\\
\bottomrule
\end{longtable}

结论是：D3 位于早期筛查阶段，应该提高召回，因此使用 hybrid；D4 位于共享写入最终边界，应该更保守，因此使用 rule-based verifier。

\section{RQ5：基线对比}

\begin{longtable}{@{}p{0.30\textwidth}p{0.10\textwidth}p{0.09\textwidth}p{0.09\textwidth}p{0.09\textwidth}p{0.09\textwidth}@{}}
\caption{与基线工具的横向对比}\\
\toprule
\textbf{Baseline} & \textbf{Status} & \textbf{Prec.} & \textbf{Rec.} & \textbf{F1} & \textbf{FPR}\\
\midrule
\endfirsthead
\toprule
\textbf{Baseline} & \textbf{Status} & \textbf{Prec.} & \textbf{Rec.} & \textbf{F1} & \textbf{FPR}\\
\midrule
\endhead
\sys{}\_ALL & ok & 1.0000 & 1.0000 & 1.0000 & 0.0000\\
PromptInject-style & ok & 1.0000 & 0.3222 & 0.4874 & 0.0000\\
LLM-Guard-style & ok & 1.0000 & 0.4333 & 0.6047 & 0.0000\\
Rebuff-style & ok & 1.0000 & 0.4222 & 0.5938 & 0.0000\\
DeepSeek-v4-flash & skipped & -- & -- & -- & --\\
\bottomrule
\end{longtable}

三类本地工具基线主要判断文本是否像提示注入，缺少共享写入网关、检索审计和执行仲裁，因此对 Multi-Agent Shared Write 攻击链覆盖不足。DeepSeek 远程基线已预留接口，但当前报告不写入 API key，默认跳过远程调用。

\section{RQ6：端到端验证}

端到端验证使用两轮任务：第一轮 \(Q_{\mathrm{inject}}\) 使 Agent A 读取恶意外部内容并尝试写入共享记忆；第二轮 \(Q_{\mathrm{target}}\) 使 Agent B 在正常任务中检索共享记忆并生成工具动作。

\begin{longtable}{@{}p{0.54\textwidth}p{0.24\textwidth}@{}}
\caption{端到端攻击链验证结果}\\
\toprule
\textbf{指标} & \textbf{结果}\\
\midrule
\endfirsthead
\toprule
\textbf{指标} & \textbf{结果}\\
\midrule
\endhead
Scenarios & 8\\
Vulnerable path dangerous executions & 8/8\\
\sys{} poisoned memory writes & 0/8\\
\sys{} dangerous executions & 0/8\\
\bottomrule
\end{longtable}

8 个场景覆盖邮件外发、数据库写入、仓库提交、生产部署和密钥读取。脆弱系统失败的原因是共享记忆采用扁平信任模型；\sys{} 在 D4 写入网关和 D6 执行仲裁两个独立边界上切断了攻击链。

\section{延迟测试}

\begin{longtable}{@{}p{0.24\textwidth}p{0.28\textwidth}p{0.10\textwidth}p{0.10\textwidth}p{0.10\textwidth}@{}}
\caption{延迟测试结果}\\
\toprule
\textbf{Experiment} & \textbf{Config} & \textbf{P50 ms} & \textbf{P95 ms} & \textbf{P99 ms}\\
\midrule
\endfirsthead
\toprule
\textbf{Experiment} & \textbf{Config} & \textbf{P50 ms} & \textbf{P95 ms} & \textbf{P99 ms}\\
\midrule
\endhead
Ablation & ALL & 0.0427 & 0.0977 & 0.1020\\
Tech Selection & ALL recommended & 0.0473 & 0.1121 & 0.1363\\
Baseline Compare & \sys{}\_ALL & 0.0485 & 0.1188 & 0.1411\\
\bottomrule
\end{longtable}

当前实现为内存版和规则/伪 rubric 版，因此延迟处于亚毫秒级。后续接入真实 LLM judge、向量库和外部工具后，需要重新评估端到端延迟和调用成本。

\chapter{总结与讨论}

\section{阶段结论}

本文围绕 Multi-Agent Shared Write 场景，给出了从攻击链到防御链的完整阶段性验证。当前结论如下：

\begin{enumerate}
  \item 共享记忆污染的本质是自动信任提升，而不是简单的单轮提示注入。
  \item 写入网关是系统不可缺少的架构边界；没有 D4，系统无法证明 shared memory 中内容的可信来源。
  \item D3 适合采用 hybrid 检测以提高隐蔽攻击召回，D4 应保持保守以控制良性同步误伤。
  \item D5 和 D6 体现纵深防御价值：污染即使进入上下文，也不能自动成为工具授权。
  \item 基础集上的指标饱和说明数据集仍需扩展，不能直接宣称系统已在真实世界完备有效。
  \item 端到端验证已经复现攻击链，并证明 \sys{} 可将危险执行从 8/8 降为 0/8。
\end{enumerate}

\section{局限性}

当前工作仍然存在局限。第一，数据集规模较小，样本主要来自公开攻击模式参考和人工重构，尚不能代表真实线上分布。第二，rubric 检测目前是确定性替身，还没有接入真实 LLM judge。第三，端到端 Agent 是最小复现系统，尚未接入 AutoGen、CrewAI 等真实多 Agent 框架。第四，当前延迟数据来自内存版实现，接入向量库、远程模型和人工审批后需要重新测量。

\section{未来工作}

下一阶段将围绕三个方向推进：一是扩展数据集，加入更多跨语言、长上下文、隐蔽同步和良性边界样本；二是补跑远程 LLM 基线，并在不泄露密钥的前提下评估成本、延迟和隐私风险；三是将 \sys{} 接入真实多 Agent 框架，验证共享记忆写入网关、检索审计和执行仲裁在真实工作流中的兼容性。

\end{document}
